\subsection{五除椭圆像的消元多项式、伽罗瓦群与切比雪夫分解密度}\label{subsec:conclusion-elliptic-five-division-chebotarev-density}
本小节在既有圆维/半圆维语义与实现层外置纪律下，建立一条独立的椭圆五除链：由重量像消元多项式出发，给出 \(\mathrm{GL}_2(\FF_5)\) 伽罗瓦刻画，并导出模素数分解型及其切比雪夫密度的全表述。

\begin{theorem}[五分点重量像的 24 次消元多项式与判别式平方核]\label{thm:conclusion-elliptic-t5-elimination-discriminant-squareclass}
设
\[
E/\QQ:\quad Y^2=X^3-X+\frac14,
\qquad
y:=X^2+Y-\frac12.
\]
定义
\[
\Pi(X,y):=X^4-X^3-(2y+1)X^2+X+y(y+1)\in\ZZ[X,y],
\]
以及五除点多项式
\[
\psi_5(X)=5X^{12}-62X^{10}+95X^9-105X^8-60X^7+285X^6-174X^5-5X^4-5X^3+35X^2-15X+2.
\]
令
\[
T_5(y):=\operatorname{Res}_X\!\bigl(\psi_5(X),\Pi(X,y)\bigr)\in\ZZ[y].
\]
则 \(T_5\) 为 24 次多项式，且
\[
\begin{aligned}
T_5(y)=\;&625 y^{24}-23500 y^{23}+244475 y^{22}+855915 y^{21}-11037859 y^{20}-89302755 y^{19}\\
&-22555640 y^{18}+761997855 y^{17}+1386419000 y^{16}-320647060 y^{15}-1928743150 y^{14}\\
&+838813180 y^{13}+3910018680 y^{12}+3220925750 y^{11}+1313468808 y^{10}-538699445 y^9\\
&-1045141370 y^8-317786725 y^7-89515935 y^6-310993491 y^5-240190705 y^4\\
&-86874340 y^3-15726875 y^2-1254110 y+10048.
\end{aligned}
\]
并满足：
\begin{enumerate}
  \item 在 \(\overline{\QQ}\) 中，
  \[
  \{\text{\(T_5\) 的根}\}
  =
  \{y(P):\ P\in E[5]\setminus\{O\}\}.
  \]
  \item \(T_5\) 在 \(\QQ[y]\) 上不可约，故对任意 \(P\in E[5]\setminus\{O\}\) 有
  \[
  [\QQ(y(P)):\QQ]=24.
  \]
  \item 存在 \(M\in\ZZ\) 使
  \[
  \operatorname{Disc}(T_5)=2^{32}\cdot 5^{35}\cdot 19^2\cdot 37^{94}\cdot M^2,
  \]
  因而 \(\operatorname{Disc}(T_5)\) 在 \(\QQ^\times/(\QQ^\times)^2\) 中的平方类为 \(5\)。
\end{enumerate}
\end{theorem}
\begin{proof}
由 \(\Pi(X,y)=0\) 可见 \(\QQ(E)=\QQ(X,Y)=\QQ(y)(X)\)，且 \(X\) 在 \(\QQ(y)\) 上满足四次关系 \(\Pi(X,y)=0\)。
若 \(P\in E[5]\setminus\{O\}\)，则 \(X(P)\) 为 \(\psi_5\) 的根，且 \(\Pi(X(P),y(P))=0\)，故
\[
T_5(y(P))=\operatorname{Res}_X(\psi_5,\Pi)(y(P))=0.
\]
反之，若 \(T_5(y_0)=0\)，由结式定义存在 \(X_0\) 使
\[
\psi_5(X_0)=0,\qquad \Pi(X_0,y_0)=0.
\]
设
\[
Y_0:=y_0-X_0^2+\frac12.
\]
由 \(\Pi\) 的构造可反推出 \((X_0,Y_0)\in E(\overline{\QQ})\)；再由 \(\psi_5(X_0)=0\) 得其为非平凡五除点，故根集刻画成立。

不可约性由模 \(3\) 约化审计给出：\(T_5\bmod 3\) 在 \(\FF_3[y]\) 上为 24 次不可约多项式，故 \(T_5\) 在 \(\QQ[y]\) 不可约（Gauss 引理）。

判别式分解由整系数判别式计算得到：提取显式非平方因子后余项为完全平方，故平方核为 \(5\)。证毕。
\end{proof}

\begin{theorem}[\(\mathrm{Gal}(T_5)\) 的 \(\mathrm{GL}_2(\FF_5)\) 刻画与全五除域识别]\label{thm:conclusion-elliptic-t5-galois-gl2f5}
设 \(L\) 为 \(T_5\) 在 \(\QQ\) 上的分裂域。则存在自然同构
\[
\mathrm{Gal}(L/\QQ)\cong \mathrm{GL}_2(\FF_5),
\]
并且
\[
L=\QQ(E[5]),
\]
从而
\[
[L:\QQ]=|\mathrm{GL}_2(\FF_5)|=(5^2-1)(5^2-5)=480.
\]
\end{theorem}
\begin{proof}
由定理 \ref{thm:conclusion-elliptic-t5-elimination-discriminant-squareclass}，\(L\) 的根置换作用可识别为对 \(E[5]\setminus\{O\}\)（24 点）的作用。
取基 \(E[5]\cong \FF_5^2\)，得到模五表示
\[
\rho_{E,5}:\mathrm{Gal}(\overline{\QQ}/\QQ)\to \mathrm{GL}_2(\FF_5),
\]
其像记为 \(G\)。于是
\[
\mathrm{Gal}(L/\QQ)\hookrightarrow G\le \mathrm{GL}_2(\FF_5),
\]
且在 24 点集合上为传递作用。

由良素数 Frobenius 循环型约束 \(G\)：
\begin{itemize}
  \item \(p=3\)：\(T_5\bmod 3\) 不可约，故 \(\mathrm{Frob}_3\) 在 24 根上为 24-循环，对应阶 24 的 Singer 元；
  \item \(p=61\)：分解型为 \((1^4\,5^4)\)，对应含有 transvection（固定一条直线 4 个非零向量，其余 20 点分成 4 个 5-循环）。
\end{itemize}
标准 Dickson 型生成结论表明：同时含有 Singer 元与 transvection 的子群必为全体 \(\mathrm{GL}_2(\FF_5)\)。
故 \(G=\mathrm{GL}_2(\FF_5)\)，从而
\[
\mathrm{Gal}(L/\QQ)\cong \mathrm{GL}_2(\FF_5).
\]
又 \(L\subseteq \QQ(E[5])\)，且两者对应同一模五 Galois 作用核，故 \(L=\QQ(E[5])\)。证毕。
\end{proof}

\begin{theorem}[\texorpdfstring{$T_5$}{T5} 分裂域的唯一二次子域恰为 \texorpdfstring{$\QQ(\sqrt5)$}{Q(sqrt5)}]\label{thm:conclusion-elliptic-t5-unique-quadratic-subfield}
\leanverified{paper\_conclusion\_elliptic\_t5\_unique\_quadratic\_subfield}
沿用定理 \ref{thm:conclusion-elliptic-t5-galois-gl2f5} 的记号，设 \(L/\QQ\) 为 \(T_5\) 的分裂域。则 \(L\) 的唯一二次子域为
\[
\QQ(\sqrt5).
\]
\end{theorem}
\begin{proof}
由定理 \ref{thm:conclusion-elliptic-t5-galois-gl2f5}，
\[
\Gal(L/\QQ)\cong \GL_2(\FF_5).
\]
二次子域与 \(\GL_2(\FF_5)\) 的指数 \(2\) 正规子群一一对应。又行列式给出满同态
\[
\GL_2(\FF_5)\twoheadrightarrow \FF_5^\times\cong C_4,
\]
并且其阿贝尔化即为 \(\FF_5^\times\)；故 \(\GL_2(\FF_5)\) 仅有一个指数 \(2\) 商，也即仅有一个指数 \(2\) 正规子群。于是 \(L\) 仅有一个二次子域。

另一方面，定理 \ref{thm:conclusion-elliptic-t5-elimination-discriminant-squareclass} 已证明 \(\Disc(T_5)\) 的平方类为 \(5\)。对不可约可分多项式，分裂域中由判别式符号角色切出的二次子域正由这一平方类决定；故该唯一二次子域只能是
\[
\QQ(\sqrt5).
\]
证毕。
\end{proof}


\begin{proposition}[\texorpdfstring{$p$}{p}-扭点像的消元多项式之满像伽罗瓦识别（\texorpdfstring{$\pm$}{±}-类版本）]\label{prop:conclusion-elliptic-ptorsion-elimination-galois-gl2-modpm}
设 \(E/\QQ\) 为椭圆曲线，\(p\) 为奇素数。
取一个有理函数 \(y:E\to\PP^1\) 满足：
\begin{enumerate}
  \item \(y(P)=y(-P)\)（即 \(y\) 在 \([-1]\) 作用下不变）；
  \item 对 \(E[p]\setminus\{O\}\) 有分离性质：\(y(P)=y(Q)\) 当且仅当 \(Q=\pm P\)。
\end{enumerate}
令 \(\psi_p(X)\in\QQ[X]\) 为 \(p\)-division 多项式的 \(X\)-坐标形式（其根集对应 \((E[p]\setminus\{O\})/\{\pm1\}\) 的 \(X\)-坐标），并取 \(\Pi_y(X,t)\in\QQ[X,t]\) 为满足 \(\Pi_y\bigl(X(P),y(P)\bigr)=0\) 的消元关系（例如将 \(y\) 写为 \(y=f(X)\in\QQ(X)\) 后清分母得到）。
定义
\[
L_{p,y}(t):=\operatorname{Res}_{X}\!\bigl(\psi_p(X),\Pi_y(X,t)\bigr)\in\QQ[t].
\]
假设模 \(p\) 伽罗瓦表示
\(
\rho_{E,p}:G_{\QQ}\to \GL_2(\FF_p)
\)
为满射。
则：
\begin{enumerate}
  \item \(L_{p,y}\) 在 \(\QQ[t]\) 上不可约；
  \item 其分裂域 \(K_{p,y}\) 的伽罗瓦群同构于
  \[
  \boxed{\ \mathrm{Gal}(K_{p,y}/\QQ)\ \cong\ \GL_2(\FF_p)/\{\pm I\}\ }
  \]
  且其对根集合的自然作用与 \(\GL_2(\FF_p)\) 对 \((\FF_p^2\setminus\{0\})/\{\pm1\}\) 的作用一致。
\end{enumerate}
\end{proposition}
\begin{proof}
由假设 (1)(2)，根集合可自然识别为 \((E[p]\setminus\{O\})/\{\pm1\}\) 在 \(y\)-坐标下的像集；因此 \(G_{\QQ}\) 在根集合上的作用由其在 \(E[p]\) 上的作用所诱导。
选定基 \(E[p]\cong \FF_p^2\) 后，\(G_{\QQ}\) 的像为 \(\GL_2(\FF_p)\)，从而在 \((\FF_p^2\setminus\{0\})/\{\pm1\}\) 上诱导出传递置换作用。
由“根轨道与不可约因子对应”可知：根的单轨道性等价于 \(L_{p,y}\) 不可约，得到 (1)。
\par
对 (2)，由于 \(\pm I\) 在 \((\FF_p^2\setminus\{0\})/\{\pm1\}\) 上作用平凡，诱导作用的核包含 \(\{\pm I\}\)。
反之，若 \(g\in\GL_2(\FF_p)\) 在全部 \(\pm\)-类上作用平凡，则对任意 \(0\neq v\in\FF_p^2\) 有 \(gv=\pm v\)。
由线性性可推出 \(g=\pm I\)，故核恰为 \(\{\pm I\}\)，从而作用群同构于 \(\GL_2(\FF_p)/\{\pm I\}\)。
该作用群即为分裂域的伽罗瓦群，结论成立。证毕。
\end{proof}

\begin{conjecture}[判别式素因子律：\texorpdfstring{$p$}{p}-扭点像消元多项式的平方类刚性]\label{conj:conclusion-elliptic-ptorsion-discriminant-squareclass-p}
在命题 \ref{prop:conclusion-elliptic-ptorsion-elimination-galois-gl2-modpm} 的设定下，进一步假设 \(E\) 在 \(p\) 处良约化。
则对除去有限个例外的奇素数 \(p\)，预期有
\[
\boxed{\ \mathrm{Disc}(L_{p,y})\ \in\ p\cdot(\QQ^\times)^2\ }.
\]
等价地，\(\QQ\bigl(\sqrt{\mathrm{Disc}(L_{p,y})}\bigr)=\QQ(\sqrt p)\)。
该平方类刚性与定理 \ref{thm:conclusion-elliptic-t5-elimination-discriminant-squareclass} 的 \(p=5\) 判别式平方核相一致，并应可解释为
\(\GL_2(\FF_p)/\{\pm I\}\) 在 \((\FF_p^2\setminus\{0\})/\{\pm1\}\) 的置换表示之符号特征与局部行列式二次特征的匹配。
\end{conjecture}

\begin{theorem}[\(T_5\) 的完整分解型与切比雪夫密度表]\label{thm:conclusion-elliptic-t5-full-chebotarev-splitting-table}
设 \(p\nmid \operatorname{Disc}(T_5)\)。
则 \(T_5\) 在 \(\FF_p[y]\) 中的不可约分解次数型仅可能为下表 14 类，且每一类的自然密度等于
\(\mathrm{GL}_2(\FF_5)\) 在 24 点作用下对应循环类型的比例（分母 \(480\)）。

\begin{table}[H]
\centering
\small
\caption{\(T_5\) 在良素数处的分解型与自然密度}
\label{tab:conclusion-elliptic-t5-factorization-density}
\begin{tabular}{lll}
\toprule
\(\FF_p\)-分解次数型 & 循环类型 & 密度 \\
\midrule
\((24)\) & \((24)\) & \(80/480=1/6\) \\
\((4,20)\) & \((4,20)\) & \(48/480=1/10\) \\
\((12,12)\) & \((12,12)\) & \(40/480=1/12\) \\
\((8,8,8)\) & \((8,8,8)\) & \(40/480=1/12\) \\
\((2,2,10,10)\) & \((2,2,10,10)\) & \(24/480=1/20\) \\
\((6,6,6,6)\) & \((6,6,6,6)\) & \(20/480=1/24\) \\
\((4,4,4,4,4,4)\) & \((4^6)\) & \(32/480=1/15\) \\
\((2,2,4,4,4,4,4)\) & \((2^2,4^5)\) & \(60/480=1/8\) \\
\((3,3,3,3,3,3,3,3)\) & \((3^8)\) & \(20/480=1/24\) \\
\((2,2,2,2,2,2,2,2,2,2,2,2)\) & \((2^{12})\) & \(1/480\) \\
\((1,1,1,1,2,2,2,2,2,2,2,2,2,2)\) & \((1^4,2^{10})\) & \(30/480=1/16\) \\
\((1,1,1,1,4,4,4,4,4)\) & \((1^4,4^5)\) & \(60/480=1/8\) \\
\((1,1,1,1,5,5,5,5)\) & \((1^4,5^4)\) & \(24/480=1/20\) \\
\((1,1,\dots,1)\)（24 个一次因子） & \((1^{24})\) & \(1/480\) \\
\bottomrule
\end{tabular}
\end{table}
\end{theorem}
\begin{proof}
由定理 \ref{thm:conclusion-elliptic-t5-galois-gl2f5}，
\[
\mathrm{Gal}(L/\QQ)\cong \mathrm{GL}_2(\FF_5)
\]
并在 24 个根上忠实作用。对任意良素数 \(p\)，\(\mathrm{Frob}_p\) 在该 24 点置换表示下的循环分解，等价于 \(T_5\bmod p\) 的不可约因子次数型。
切比雪夫密度定理给出：每种分解型的自然密度等于对应共轭类（并类并）在群中的比例。对 \(\mathrm{GL}_2(\FF_5)\) 的 24 点作用逐类计数即得上表。证毕。
\end{proof}

\begin{corollary}[模素数可解性与完全分裂密度]\label{cor:conclusion-elliptic-t5-solvability-and-total-splitting-density}
\leanverified{paper\_conclusion\_elliptic\_t5\_solvability\_and\_total\_splitting\_density}
在定理 \ref{thm:conclusion-elliptic-t5-full-chebotarev-splitting-table} 的设定下：
\begin{enumerate}
  \item \(T_5\) 在 \(\FF_p\) 上存在一次因子的素数 \(p\) 的自然密度为
  \[
  \frac{23}{96}.
  \]
  \item \(T_5\) 在 \(\FF_p\) 上完全分裂（24 个一次因子）的素数 \(p\) 的自然密度为
  \[
  \frac{1}{480}.
  \]
\end{enumerate}
\end{corollary}
\begin{proof}
含一次因子的分解型为
\[
(1^4,2^{10}),\ (1^4,4^5),\ (1^4,5^4),\ (1^{24}).
\]
由表 \ref{tab:conclusion-elliptic-t5-factorization-density} 求和：
\[
\frac{30}{480}+\frac{60}{480}+\frac{24}{480}+\frac{1}{480}
=\frac{115}{480}
=\frac{23}{96}.
\]
完全分裂对应恒等元共轭类，密度 \(1/|\mathrm{GL}_2(\FF_5)|=1/480\)。证毕。
\end{proof}

\begin{theorem}[\texorpdfstring{$T_5$}{T5} 的线性因子数在模 \texorpdfstring{$5$}{5} 剩余类上的条件分布与统一条件期望]\label{thm:conclusion-elliptic-t5-linear-factor-conditional-density-mod5}
沿用定理 \ref{thm:conclusion-elliptic-t5-galois-gl2f5} 的记号。对每个良素数
\[
\ell\nmid \operatorname{Disc}(T_5),
\]
记
\[
N_\ell:=\#\{\text{\(T_5\bmod \ell\) 的一次因子}\}.
\]
则按 \(\ell\bmod 5\) 条件化之后，\(N_\ell\) 的分布满足：
\begin{enumerate}
  \item 当 \(\ell\equiv 1\pmod 5\) 时，
  \[
  \delta\!\bigl(N_\ell=24\mid \ell\equiv 1\!\!\!\!\pmod 5\bigr)=\frac{1}{120},
  \]
  \[
  \delta\!\bigl(N_\ell=4\mid \ell\equiv 1\!\!\!\!\pmod 5\bigr)=\frac{24}{120}=\frac15,
  \]
  \[
  \delta\!\bigl(N_\ell=0\mid \ell\equiv 1\!\!\!\!\pmod 5\bigr)=\frac{95}{120}=\frac{19}{24}.
  \]
  \item 对每个 \(a\in\{2,3,4\}\)，都有
  \[
  \delta\!\bigl(N_\ell=4\mid \ell\equiv a\!\!\!\!\pmod 5\bigr)=\frac14,
  \qquad
  \delta\!\bigl(N_\ell=0\mid \ell\equiv a\!\!\!\!\pmod 5\bigr)=\frac34.
  \]
\end{enumerate}
此外，对每个 \(a\in(\ZZ/5\ZZ)^\times\) 都有
\[
\mathbb E\!\bigl(N_\ell\mid \ell\equiv a\!\!\!\!\pmod 5\bigr)=1.
\]
特别地，\(\ell\equiv 1\pmod 5\) 是 \(T_5\bmod \ell\) 完全分裂的唯一可能同余类；而当 \(\ell\not\equiv 1\pmod 5\) 时，\(T_5\bmod \ell\) 不可能完全分裂，且“一次因子存在”这一事件的条件密度恒为 \(1/4\)。
\end{theorem}
\begin{proof}
由定理 \ref{thm:conclusion-elliptic-t5-galois-gl2f5}，\(\Gal(L/\QQ)\cong \GL_2(\FF_5)\)，且其在 \(T_5\) 的 \(24\) 个根上的作用可识别为 \(\GL_2(\FF_5)\) 在
\[
V^\times:=\FF_5^2\setminus\{0\}
\]
上的自然作用。故对每个良素数 \(\ell\)，Frobenius 元 \(g_\ell:=\Frob_\ell\) 在 \(V^\times\) 上的不动点数正好等于 \(N_\ell\)。

另一方面，由 Weil pairing 的标准性质，
\[
\det(g_\ell)\equiv \ell \pmod 5.
\]
因此，对每个 \(a\in\FF_5^\times\)，条件 \(\ell\equiv a\pmod 5\) 等价于把 Frobenius 限制在 determinant 纤维
\[
G^{(a)}:=\{g\in \GL_2(\FF_5):\ \det g=a\}.
\]
由于 \(\det:\GL_2(\FF_5)\to \FF_5^\times\) 为满射，故
\[
|G^{(a)}|=\frac{|\GL_2(\FF_5)|}{|\FF_5^\times|}=\frac{480}{4}=120.
\]

先设 \(a\neq 1\)。若 \(g\in G^{(a)}\) 在 \(V^\times\) 上有不动点，则 \(1\) 是其特征值；又因 \(\det g=a\neq 1\)，另一特征值必为 \(a\)，故 \(g\) 与 \(\diag(1,a)\) 共轭。此时中央化子由同一特征基下的全部可逆对角矩阵组成，大小为
\[
(5-1)^2=16,
\]
故该共轭类大小为
\[
\frac{480}{16}=30.
\]
这些矩阵恰固定其 \(1\)-特征线上的 \(4\) 个非零向量，因此对应 \(N_\ell=4\)。其余
\[
120-30=90
\]
个 determinant 为 \(a\) 的矩阵无非零不动点，对应 \(N_\ell=0\)。于是
\[
\delta\!\bigl(N_\ell=4\mid \ell\equiv a\!\!\!\!\pmod 5\bigr)=\frac{30}{120}=\frac14,
\qquad
\delta\!\bigl(N_\ell=0\mid \ell\equiv a\!\!\!\!\pmod 5\bigr)=\frac{90}{120}=\frac34.
\]

再设 \(a=1\)。此时 \(G^{(1)}=\mathrm{SL}_2(\FF_5)\) 共有 \(120\) 个元素。恒等元固定全部 \(24\) 个非零向量，故对应 \(N_\ell=24\)，其条件密度为 \(1/120\)。若 \(g\neq I\) 且在 \(V^\times\) 上有非零不动点，则其特征多项式只能为 \((X-1)^2\)，故 \(g\) 是非平凡幺幂元，与
\[
J=\begin{pmatrix}1&1\\0&1\end{pmatrix}
\]
共轭。其中央化子为
\[
\{aI+b(J-I):\ a\in \FF_5^\times,\ b\in \FF_5\},
\]
大小为 \(4\cdot 5=20\)，故该共轭类大小为
\[
\frac{480}{20}=24.
\]
每个此类矩阵的 \(1\)-特征空间是一条直线，因而恰固定 \(4\) 个非零向量，即对应 \(N_\ell=4\)。其余
\[
120-1-24=95
\]
个 determinant 为 \(1\) 的矩阵无非零不动点，对应 \(N_\ell=0\)。遂得所述三项条件密度。

最后计算条件期望。若 \(a\neq 1\)，则
\[
\mathbb E\!\bigl(N_\ell\mid \ell\equiv a\!\!\!\!\pmod 5\bigr)
=
\frac{30\cdot 4}{120}
=
1.
\]
若 \(a=1\)，则
\[
\mathbb E\!\bigl(N_\ell\mid \ell\equiv 1\!\!\!\!\pmod 5\bigr)
=
\frac{1\cdot 24+24\cdot 4}{120}
=
1.
\]
故对全部非零剩余类都有
\[
\mathbb E\!\bigl(N_\ell\mid \ell\equiv a\!\!\!\!\pmod 5\bigr)=1.
\]
证毕。
\end{proof}

\begin{theorem}[\texorpdfstring{$T_5\bmod \ell$}{T5 mod ell} 的不可约因子个数满足二次特征奇偶律]\label{thm:conclusion-elliptic-t5-irreducible-factor-parity-legendre}
沿用定理 \ref{thm:conclusion-elliptic-t5-galois-gl2f5} 的记号。对每个良素数 \(\ell\nmid \operatorname{Disc}(T_5)\)，把
\[
T_5(y)\bmod \ell=\prod_{j=1}^{r_\ell}P_j(y)
\]
写成 \(\FF_\ell[y]\) 中的互异不可约因子分解，并记 \(r_\ell\) 为不可约因子个数。则有严格奇偶律
\[
(-1)^{r_\ell}=\left(\frac{\ell}{5}\right),
\]
其中 \(\left(\frac{\ell}{5}\right)\) 为 Legendre 符号。等价地，
\[
\ell\equiv 1,4\!\!\!\!\pmod 5 \Longrightarrow r_\ell\equiv 0\!\!\!\!\pmod 2,
\qquad
\ell\equiv 2,3\!\!\!\!\pmod 5 \Longrightarrow r_\ell\equiv 1\!\!\!\!\pmod 2.
\]
\end{theorem}
\begin{proof}
由定理 \ref{thm:conclusion-elliptic-t5-galois-gl2f5}，\(T_5\) 的根集合可识别为 \(V^\times=\FF_5^2\setminus\{0\}\)，而良素数 \(\ell\) 的 Frobenius 元 \(g_\ell\) 在该 \(24\) 点集上的循环分解与 \(T_5\bmod \ell\) 的不可约分解型一致。因此若循环个数为 \(r_\ell\)，则其置换符号满足
\[
\operatorname{sgn}(g_\ell)=(-1)^{24-r_\ell}=(-1)^{r_\ell},
\]
因为 \(24\) 为偶数。

现定义
\[
\sigma:\GL_2(\FF_5)\longrightarrow \{\pm 1\},
\qquad
\sigma(g):=\operatorname{sgn}(g\curvearrowright V^\times).
\]
这是一个群同态。先看幺幂元
\[
J=\begin{pmatrix}1&1\\0&1\end{pmatrix}.
\]
它固定直线 \(\FF_5 e_1\) 上的 \(4\) 个非零向量，而把其余 \(20\) 个向量分成 \(4\) 个 \(5\)-循环，故
\[
\sigma(J)=1.
\]
因此 \(\sigma\) 在 \(\mathrm{SL}_2(\FF_5)\) 的标准幺幂生成元上取值为 \(1\)，从而在整个 \(\mathrm{SL}_2(\FF_5)\) 上恒等于 \(1\)。

再取对角元
\[
D_a:=\diag(a,1)\qquad (a\in \FF_5^\times).
\]
在集合
\[
H_0:=\{(x,0):\ x\in \FF_5^\times\}
\]
上，\(D_a\) 的作用即为 \(\FF_5^\times\) 上的乘法 \(x\mapsto ax\)，其置换符号在 \(a=1,4\) 时为 \(+1\)，在 \(a=2,3\) 时为 \(-1\)，恰等于 \(\bigl(\frac{a}{5}\bigr)\)。
对每个 \(b\in \FF_5^\times\)，在
\[
H_b:=\{(x,b):\ x\in \FF_5\}
\]
上，\(D_a\) 固定 \((0,b)\)，并在其余 \(4\) 个点上施行同一乘法置换，故其符号仍为 \(\bigl(\frac{a}{5}\bigr)\)。由于这样的 \(H_b\) 有 \(4\) 个，它们贡献的总符号为
\[
\left(\frac{a}{5}\right)^4=1.
\]
因此
\[
\sigma(D_a)=\left(\frac{a}{5}\right).
\]

任取 \(g\in \GL_2(\FF_5)\)。把它写成 \(g=sh\)，其中 \(s\in \mathrm{SL}_2(\FF_5)\)、\(h=D_{\det g}\)。于是
\[
\sigma(g)=\sigma(s)\sigma(h)=\left(\frac{\det g}{5}\right).
\]
对 \(g=g_\ell\) 应用此式，并利用 Weil pairing 给出的
\[
\det g_\ell\equiv \ell \pmod 5,
\]
即可得到
\[
\operatorname{sgn}(g_\ell)=\left(\frac{\ell}{5}\right).
\]
与前述 \(\operatorname{sgn}(g_\ell)=(-1)^{r_\ell}\) 比较，结论成立。证毕。
\end{proof}

\begin{corollary}[\texorpdfstring{$T_5$}{T5} 的分解奇偶谱唯一恢复二次分辨子域]\label{cor:conclusion-elliptic-t5-parity-spectrum-recovers-qsqrt5}
沿用定理 \ref{thm:conclusion-elliptic-t5-irreducible-factor-parity-legendre} 的记号。定义良素数集合
\[
\mathcal P_{\mathrm{odd}}
:=
\{\ell\nmid \Disc(T_5):\ r_\ell\ \text{为奇数}\}.
\]
则
\[
\boxed{\;
\mathcal P_{\mathrm{odd}}
=
\{\ell\nmid \Disc(T_5):\ \ell\ \text{在 }\QQ(\sqrt5)\text{ 中惰性}\}. \;}
\]
从而：
\begin{enumerate}
  \item \(\mathcal P_{\mathrm{odd}}\) 的自然密度等于 \(1/2\)；
  \item 若 \(K/\QQ\) 为任意二次域，且除有限个素数外有
        \[
        \mathbf 1_{\ell\text{ inert in }K}
        =
        \mathbf 1_{\ell\in\mathcal P_{\mathrm{odd}}},
        \]
        则必有
        \[
        K=\QQ(\sqrt5).
        \]
\end{enumerate}
换言之，仅由 \(T_5\bmod \ell\) 的不可约因子总数奇偶性，便可唯一恢复其二次分辨子域。
\end{corollary}
\begin{proof}
由定理 \ref{thm:conclusion-elliptic-t5-irreducible-factor-parity-legendre}，
\[
\ell\in\mathcal P_{\mathrm{odd}}
\iff
(-1)^{r_\ell}=-1
\iff
\left(\frac{\ell}{5}\right)=-1.
\]
对奇素数 \(\ell\neq 5\)，由于 \(5\equiv 1\pmod 4\)，二次互反律给出
\[
\left(\frac{\ell}{5}\right)=\left(\frac{5}{\ell}\right).
\]
而对二次域 \(\QQ(\sqrt5)\)，除坏素数外，\(\ell\) 惰性当且仅当
\[
\left(\frac{5}{\ell}\right)=-1.
\]
故所述集合恒等式成立。

第一条随即由非平凡二次 Dirichlet 特征取值 \(\pm1\) 的 Chebotarev/Dirichlet 密度各占一半得到。

第二条中，设 \(K=\QQ(\sqrt d)\) 为某个二次域，并假设除有限个素数外有
\[
\mathbf 1_{\ell\text{ inert in }K}
=
\mathbf 1_{\ell\in\mathcal P_{\mathrm{odd}}}.
\]
则对应二次特征
\[
\chi_d(\ell):=\left(\frac{d}{\ell}\right)
\]
与
\[
\chi_5(\ell):=\left(\frac{5}{\ell}\right)
\]
在几乎所有素数上相等。由二次 Dirichlet 特征的唯一性，必有
\[
\chi_d=\chi_5,
\]
从而
\[
K=\QQ(\sqrt5).
\]
这也与定理 \ref{thm:conclusion-elliptic-t5-unique-quadratic-subfield} 中 \(T_5\) 分裂域的唯一二次子域结论完全一致。证毕。
\end{proof}

\begin{proposition}[\(T_5\) 的实根结构与生成域签名]\label{prop:conclusion-elliptic-t5-real-root-signature}
\(T_5(y)\) 在 \(\RR\) 上恰有 4 个互异实根；因此对任意根 \(\theta\)，数域
\[
K_\theta:=\QQ(\theta)
\]
的签名为
\[
(r_1,r_2)=(4,10).
\]
并且四个实根可由严格分离区间定位于
\[
(0,0.01),\quad (0.7,0.8),\quad (3.1,3.2),\quad (12.2,12.3),
\]
对应数值近似
\[
0.0073136,\quad 0.7764323,\quad 3.1132063,\quad 12.2187073.
\]
\end{proposition}
\begin{proof}
对 \(T_5\) 构造 Sturm 链并比较 \(\pm\infty\) 与各区间端点的变号数，可得实根总数为 4 且均为单根。
再由区间端点符号与 Sturm 计数，排除负实根并确定上述四个分离区间。
由于 \([K_\theta:\QQ]=24\)，且实嵌入个数为 4，故
\[
24=r_1+2r_2,\qquad r_1=4
\]
解得 \(r_2=10\)。证毕。
\end{proof}

\begin{theorem}[\(\mathrm{GL}_2(\FF_q)\) 的迹层计数]\label{thm:conclusion-gl2fq-trace-strata-count}
设 \(q\) 为奇素数幂，并记
\[
G:=\mathrm{GL}_2(\FF_q).
\]
对每个 \(t\in\FF_q\)，令
\[
N_q(t):=\#\{g\in G:\ \mathrm{tr}(g)=t\}.
\]
则
\[
N_q(0)=q^3-q^2,
\qquad
N_q(t)=q^3-q^2-q\quad(t\neq 0).
\]
从而迹为 \(0\) 的比例与每个非零迹层的比例分别为
\[
\frac{q^3-q^2}{(q^2-1)(q^2-q)},
\qquad
\frac{q^3-q^2-q}{(q^2-1)(q^2-q)}.
\]
\end{theorem}
\begin{proof}
固定 \(t\in\FF_q\)。
满足 \(\mathrm{tr}(g)=t\) 的全部 \(2\times 2\) 矩阵可写为
\[
g=
\begin{pmatrix}
a & b\\
c & t-a
\end{pmatrix},
\qquad
a,b,c\in\FF_q,
\]
故总数为 \(q^3\)。
其中奇异条件为
\[
\det(g)=a(t-a)-bc=0.
\]
对固定 \(a\)，记
\[
u(a):=a(t-a).
\]
若 \(u(a)=0\)，则方程 \(bc=0\) 有 \(2q-1\) 组解；若 \(u(a)\neq 0\)，则对每个 \(b\neq 0\) 唯一决定 \(c=u(a)/b\)，故共有 \(q-1\) 组解。

当 \(t=0\) 时，\(u(a)=-a^2\) 仅在 \(a=0\) 处为零，因此奇异矩阵总数为
\[
(2q-1)+(q-1)(q-1)=q^2.
\]
故
\[
N_q(0)=q^3-q^2.
\]

当 \(t\neq 0\) 时，\(u(a)=a(t-a)\) 在 \(a=0,t\) 两处为零，因此奇异矩阵总数为
\[
2(2q-1)+(q-2)(q-1)=q^2+q.
\]
故
\[
N_q(t)=q^3-(q^2+q)=q^3-q^2-q
\qquad (t\neq 0).
\]
最后除以
\[
|G|=(q^2-1)(q^2-q)
\]
即得密度公式。证毕。
\end{proof}
\leanverified{paper\_conclusion\_gl2fq\_trace\_strata\_count}

\begin{corollary}[\(T_5\) 分裂域上的模 \(5\) Frobenius 迹密度]\label{cor:conclusion-elliptic-t5-frobenius-trace-density-mod5}
沿用定理 \ref{thm:conclusion-elliptic-t5-galois-gl2f5} 的记号，设 \(L/\QQ\) 为 \(T_5\) 的分裂域。
对每个在 \(L\) 中不分歧的素数 \(p\)，把 Frobenius 共轭类视为
\[
\mathrm{Frob}_p\in \mathrm{Gal}(L/\QQ)\cong \mathrm{GL}_2(\FF_5).
\]
则
\[
\delta\bigl(\mathrm{tr}(\mathrm{Frob}_p)\equiv 0\!\!\pmod 5\bigr)=\frac{100}{480}=\frac{5}{24},
\]
并且对每个 \(t\in\{1,2,3,4\}\) 有
\[
\delta\bigl(\mathrm{tr}(\mathrm{Frob}_p)\equiv t\!\!\pmod 5\bigr)=\frac{95}{480}=\frac{19}{96}.
\]
\end{corollary}
\begin{proof}
由定理 \ref{thm:conclusion-elliptic-t5-galois-gl2f5}，
\[
\mathrm{Gal}(L/\QQ)\cong \mathrm{GL}_2(\FF_5),
\qquad
|\mathrm{GL}_2(\FF_5)|=(25-1)(25-5)=480.
\]
再由定理 \ref{thm:conclusion-gl2fq-trace-strata-count} 在 \(q=5\) 时得到
\[
\#\mathrm{tr}^{-1}(0)=5^3-5^2=100,
\qquad
\#\mathrm{tr}^{-1}(t)=5^3-5^2-5=95\quad(t\neq 0).
\]
Chebotarev 密度定理给出：任一共轭不变子集的自然密度等于其群比例。
而迹层是共轭不变的，故结论成立。证毕。
\end{proof}

\begin{conjecture}[Lee--Yang 边界测度与 \(\mathrm{GL}_2\) Satake 参数的代数推送对偶]\label{conj:conclusion-leyang-gl2-satake-boundary-duality}
沿用定理 \ref{thm:xi-terminal-zm-leyang-elliptic-structure} 的 Lee--Yang 谱曲线
\[
\Pi(\lambda,y)=0
\]
与定理 \ref{thm:xi-terminal-zm-leyang-elliptic-modularity-37} 的椭圆模性锚点 \(E_0/\QQ\)。
设对每个良素数 \(p\nmid 37\)，\((\alpha_p,\beta_p)\) 为 \(E_0\) 的 Satake 参数，满足
\[
\alpha_p\beta_p=p,
\qquad
|\alpha_p|=|\beta_p|=p^{1/2}.
\]
则预期存在由 \(\Pi(\lambda,y)\) 的分支数据决定的显式代数映射
\[
\Phi:\TT_{\mathrm{Sat}}\to \TT,
\qquad
\TT_{\mathrm{Sat}}:=\{(u,u^{-1}):u\in\TT\},
\]
使得经验测度
\[
\nu_X
:=
\frac{1}{\#\{p\le X:\ p\nmid 37\}}
\sum_{\substack{p\le X\\ p\nmid 37}}
\delta_{\Phi(\alpha_p/p^{1/2},\,\beta_p/p^{1/2})}
\]
在 \(X\to\infty\) 时弱收敛到一个由 Lee--Yang 边界决定的概率测度 \(\mu_{\mathrm{LY}}\)。

并且，这一收敛不应具有统一可计算的 \(W_2\)-误差模。
事实上，在紧度量空间上有 \(W_1\le W_2\)；故若存在统一可计算的 \(W_2\) 控制，则会推出统一可计算的 \(W_1\) 控制。
因此，在把 \(\mu_{\mathrm{LY}}\) 识别为附录定理 \ref{thm:fold-computability-halting-no-uniform-w1-learning} 所构造的 Lee--Yang 极限测度的任何实现中，该统一模都与既有的 \(W_1\) 不可统一逼近障碍不相容。

在这一意义下，推论 \ref{cor:conclusion-elliptic-t5-frobenius-trace-density-mod5} 所给出的模 \(5\) 迹密度可视为该推送对偶在有限域层上的离散阴影。
\end{conjecture}
